\documentclass[11pt]{article}
\usepackage[T1]{fontenc}
\usepackage{amsmath,amssymb,amsthm}
\usepackage[margin=1in]{geometry}
\usepackage[hidelinks]{hyperref}
\hypersetup{
  pdftitle={A Modulo-Five Congruence for Overcolored Partition 4-Tuples Restricted by Parity of the Parts},
  pdfauthor={Xiaojian Zeng},
  pdfsubject={Partition congruences and modular forms},
  pdfkeywords={overcolored partitions, congruences, eta quotients, Sturm bound}
}

\newtheorem{theorem}{Theorem}
\newtheorem{lemma}[theorem]{Lemma}
\newcommand{\Ffive}{\mathbb F_5}
\newcommand{\ord}{\operatorname{ord}}

\title{A Modulo-Five Congruence for Overcolored Partition\
       $4$-Tuples Restricted by Parity of the Parts}
\author{Xiaojian Zeng\\
PBC School of Finance, Tsinghua University, 100190\\
\texttt{xiaojianzeng@foxmail.com}}
\date{}

\begin{document}
\maketitle

\begin{abstract}
Let $f_j=(q^j;q^j)_\infty$ and define $a(N)$ by
\[
 \frac{f_4^8}{f_1^{16}f_2^8}=\sum_{N\geq0}a(N)q^N.
\]
We prove that $a(5n)\equiv0\pmod 5$ for every $n\geq1$, resolving the
first displayed congruence in Conjecture~7.1 of Thejitha and Fathima.
The proof places a holomorphic eta quotient in
$M_{52}(\Gamma_0(20))$, applies $U_5$, and uses Sturm's theorem to reduce
the required identity to 181 exact coefficient comparisons over
$\mathbb F_5$.
\end{abstract}

\section{Statement and context}

For $j\geq1$, put
\[
 f_j=(q^j;q^j)_\infty=\prod_{m\geq1}(1-q^{jm}).
\]
The overcolored partition $k$-tuples restricted by parity studied by
Thejitha and Fathima have generating function
\begin{equation}\label{eq:general-gf}
 \sum_{N\geq0}\bar b^k_{r,s}(N)q^N
 =\frac{f_2^{(3s-2r)k}}{f_1^{2sk}f_4^{(s-r)k}}.
\end{equation}
Their Conjecture~7.1 records several congruences supported there by
computation~\cite{ThejithaFathima2026}.  The specialization
$(r,s,k)=(4,2,4)$ in \eqref{eq:general-gf} gives
\begin{equation}\label{eq:F}
 F(q):=\frac{f_4^8}{f_1^{16}f_2^8}
      =\sum_{N\geq0}a(N)q^N,
 \qquad a(N)=\bar b^4_{4,2}(N).
\end{equation}
Related generating-function congruences for the overcubic specialization
were obtained by Chac\'on and Sellers~\cite{ChaconSellers2026}; their
families do not contain the specialization in \eqref{eq:F}.

\begin{theorem}\label{thm:main}
For every integer $n\geq1$,
\[
 \bar b^4_{4,2}(5n)\equiv0\pmod5.
\]
\end{theorem}

The theorem concerns all $n\geq1$; a finite coefficient search alone would
not prove it.  We use the finite comparison only after a modular-form
reduction makes Sturm's bound applicable.

\section{Eta-quotient reduction}

Write $q=e^{2\pi iz}$ and let
\[
 U_5\!\left(\sum_{N\geq0}c(N)q^N\right)
   =\sum_{n\geq0}c(5n)q^n.
\]
Since $a(0)=1$, Theorem~\ref{thm:main} is equivalent to
\begin{equation}\label{eq:target}
 U_5(F)\equiv1\pmod5.
\end{equation}

Consider
\begin{equation}\label{eq:A}
 A(z)=\frac{\eta(4z)^8\eta(5z)^{120}}
             {\eta(z)^{16}\eta(2z)^8}.
\end{equation}
At level $20$, its nonzero eta exponents are
\[
 (r_1,r_2,r_4,r_5)=(-16,-8,8,120).
\]
The eta-quotient modularity conditions originating with Newman, together
with the cusp-order criterion of Ligozat, will be used in their standard
form~\cite{Newman1959,Ligozat1975}.  Here the weight is
$\frac12\sum r_\delta=52$, and
\[
 \sum_{\delta\mid20}\delta r_\delta=600\equiv0\pmod{24},
 \qquad
 \sum_{\delta\mid20}\frac{20}{\delta}r_\delta
 =120\equiv0\pmod{24}.
\]
Moreover,
\[
 \prod_{\delta\mid20}\delta^{r_\delta}=2^8 5^{120}
\]
is a square, so the associated quadratic character is trivial.  For a
cusp of denominator $d\mid20$, the invariant order is
\begin{equation}\label{eq:cusp-order}
 v_d(A)=\frac{20}{24\gcd(d,20/d)}
 \sum_{\delta\mid20}
 \frac{\gcd(d,\delta)^2r_\delta}{d\delta}.
\end{equation}
Direct substitution gives
\[
\begin{array}{c|rrrrrr}
 d&1&2&4&5&10&20\\ \hline
 v_d(A)&5&0&5&97&24&25.
\end{array}
\]
All cusp orders are nonnegative.  Consequently
\begin{equation}\label{eq:modularity-A}
 A\in M_{52}(\Gamma_0(20))
\end{equation}
with trivial character.  Since $5\mid20$, the standard $U_5$ operator
preserves this space~\cite{AtkinLehner1970}; hence
$U_5(A)\in M_{52}(\Gamma_0(20))$.

\section{Sturm comparison and proof}

Let
\[
 E_4(z)=1+240\sum_{n\geq1}\sigma_3(n)q^n,
 \qquad
 \Delta(z)=\eta(z)^{24}=qf_1^{24},
\]
and define
\begin{equation}\label{eq:C}
 C(z)=E_4(z)^2U_5(A)(z)-\Delta(z)^5.
\end{equation}
Both summands have weight $60$, level dividing $20$, trivial character,
and integral Fourier coefficients.  Thus
\[
 C\in M_{60}(\Gamma_0(20))\cap\mathbb Z[[q]].
\]
The index and the resulting Sturm bound are
\begin{equation}\label{eq:sturm-bound}
 [\mathrm{SL}_2(\mathbb Z):\Gamma_0(20)]
 =20\left(1+\frac12\right)\left(1+\frac15\right)=36,
 \qquad
 \frac{60}{12}\,36=180.
\end{equation}
Sturm's congruence theorem~\cite{Sturm1987} therefore shows that
$C\equiv0\pmod5$ once its coefficients through degree $180$ vanish modulo
$5$.

This finite check was carried out directly in $\Ffive[[q]]$.  Expanding
the defining Euler products through degree $900=5\cdot180$ gives
\[
 [q^{5n}]A=[q^n]\Delta^5,
 \qquad 0\leq n\leq180.
\]
For completeness, the accompanying verifier starts with the appropriate
monomial shift and processes each factor $(1-q^{jm})^e$ by coefficient
updates: for $e>0$ it performs $e$ descending subtraction passes, while
for $e<0$ it performs $-e$ ascending addition passes.  All entries are
reduced modulo $5$.  It constructs $A$ through degree $900$ and
$\Delta^5=q^5f_1^{120}$ through degree $180$, then compares precisely the
181 displayed pairs.  The failure list is empty.  Since
$E_4^2\equiv1\pmod5$, these comparisons are exactly the coefficients of
\eqref{eq:C} required by \eqref{eq:sturm-bound}.  Hence
\begin{equation}\label{eq:Czero}
 C\equiv0\pmod5.
\end{equation}

It remains to recover the original series.  Put
$D(q)=q^5f_1^{120}=\Delta(q)^5$.  The eta shifts give the exact
formal-power-series identity
\begin{equation}\label{eq:factorization}
 A(q)=q^{25}f_5^{120}F(q)=D(q^5)F(q).
\end{equation}
For arbitrary formal series $D$ and $G$, coefficient convolution yields
\begin{equation}\label{eq:u-factorization}
 U_5\bigl(D(q^5)G(q)\bigr)=D(q)U_5(G(q)).
\end{equation}
Equations \eqref{eq:factorization} and \eqref{eq:u-factorization} imply
\[
 U_5(A)=D(q)U_5(F)=\Delta^5U_5(F).
\]
Reducing \eqref{eq:Czero} modulo $5$ and again using
$E_4^2\equiv1\pmod5$, we obtain
\[
 0\equiv\Delta^5\bigl(U_5(F)-1\bigr)\pmod5.
\]
The series $\Delta^5=q^5f_1^{120}$ is nonzero in the integral domain
$\Ffive[[q]]$ (indeed, $f_1^{120}$ is a unit), so it can be cancelled.
This proves \eqref{eq:target}, and hence Theorem~\ref{thm:main}.

\section*{Computational reproducibility and assistance}

The accompanying file \texttt{modular\_sturm\_certificate.py} uses only
exact integer arithmetic reduced modulo $5$ and requires no random seed or
numerical tolerance.  Its role is limited to the Sturm-prescribed finite
comparison above.  The broader practice of separating a mathematical
reduction from a small deterministic verifier also appears in an unrelated
finite-geometric setting in Zeng et al.~\cite{ZengEtAl2026}; the present
certificate uses Euler-product arithmetic rather than enumeration or
optimization.

OpenAI Codex assisted with literature search, proof organization, exact
coefficient verification, and manuscript preparation.  Language-model
output is not used as mathematical evidence: the proof is given above and
the finite step is reproducible by the exact verifier.  No proof assistant
was used.

\begingroup
\raggedright
\bibliographystyle{plain}
\bibliography{references}
\endgroup

\end{document}
